\editionPage[\pdfbookmark[0]{Défis d’approfondissement}{bm-defis-debut-parent}]{defis-debut}{Dix-huit défis}{Un défi par chapitre · Parcours facultatif}{Chercher avant de consulter les deux indications.}
\begin{blocEdition}{defis-debut}
Ces problèmes demandent davantage d’initiative. Ils utilisent les outils du livre ; les prolongements sont définis dans l’énoncé. Les défis sont facultatifs et approfondissent les outils du programme.\par\bigskip
\begin{savoirfaire}[Une recherche qui avance]Faire un schéma, tester un petit cas, comparer deux méthodes, conserver les essais et préciser ce qui reste à prouver. Ouvrir la première indication seulement après une recherche réelle ; revenir au problème avant de lire la seconde.\end{savoirfaire}
\begin{savoirfaire}[Après la solution]Identifier l’idée décisive et la refaire sur une feuille vierge. Expliquer pourquoi elle fonctionne compte autant que retrouver le résultat.\end{savoirfaire}
\end{blocEdition}
\editionPage{defi-ex-1}{D01 · Défi}{Cent termes, sans cent calculs}{Approfondissement · 30 à 60 minutes indicatives}
\begin{blocEdition}{defi-ex-1}
Calculer $S=-1+2-3+4-\cdots-99+100$ sans additionner terme par terme. Expliquer le résultat. Que vaut la somme si l’on s’arrête à $-99$ ?\par\vspace{5mm}{\otLight\fontsize{8}{10}\selectfont\color{otMuted}Espace de recherche}\par\par\vspace{7mm}{\color{otRule!65}\hrule height .25pt}\par\vspace{7mm}{\color{otRule!65}\hrule height .25pt}\par\vspace{7mm}{\color{otRule!65}\hrule height .25pt}\par\vspace{7mm}{\color{otRule!65}\hrule height .25pt}\par\vspace{7mm}{\color{otRule!65}\hrule height .25pt}\par\vspace{7mm}{\color{otRule!65}\hrule height .25pt}
\end{blocEdition}
\editionPage{defi-ex-2}{D02 · Défi}{Retrouver tout le budget}{Approfondissement · 30 à 60 minutes indicatives}
\begin{blocEdition}{defi-ex-2}
On utilise un tiers d’un budget, puis un quart du reste, puis trois cinquièmes du nouveau reste. Il reste alors 18 euros. Retrouver le budget et les trois sommes utilisées.\par\vspace{5mm}{\otLight\fontsize{8}{10}\selectfont\color{otMuted}Espace de recherche}\par\par\vspace{7mm}{\color{otRule!65}\hrule height .25pt}\par\vspace{7mm}{\color{otRule!65}\hrule height .25pt}\par\vspace{7mm}{\color{otRule!65}\hrule height .25pt}\par\vspace{7mm}{\color{otRule!65}\hrule height .25pt}\par\vspace{7mm}{\color{otRule!65}\hrule height .25pt}\par\vspace{7mm}{\color{otRule!65}\hrule height .25pt}
\end{blocEdition}
\editionPage{defi-ex-3}{D03 · Défi}{Jusqu’où peut-on doubler ?}{Approfondissement · 30 à 60 minutes indicatives}
\begin{blocEdition}{defi-ex-3}
Une mémoire fictive contient au départ un seul élément. À chaque étape, le nombre d’éléments double. Sa capacité maximale est d’un million d’éléments. Quel est le plus grand nombre d’étapes possible ? Donner un encadrement d’un million par deux puissances de 2 consécutives.\par\vspace{5mm}{\otLight\fontsize{8}{10}\selectfont\color{otMuted}Espace de recherche}\par\par\vspace{7mm}{\color{otRule!65}\hrule height .25pt}\par\vspace{7mm}{\color{otRule!65}\hrule height .25pt}\par\vspace{7mm}{\color{otRule!65}\hrule height .25pt}\par\vspace{7mm}{\color{otRule!65}\hrule height .25pt}\par\vspace{7mm}{\color{otRule!65}\hrule height .25pt}\par\vspace{7mm}{\color{otRule!65}\hrule height .25pt}
\end{blocEdition}
\editionPage{defi-ex-4}{D04 · Défi}{Racine de deux, un nombre non décimal}{Approfondissement · 30 à 60 minutes indicatives}
\begin{blocEdition}{defi-ex-4}
Le carré construit sur la diagonale d’un carré de côté 1 a une aire égale à 2 : cette diagonale mesure $\sqrt2$. On veut démontrer par l’absurde que $\sqrt2$ n’est pas un nombre décimal.\begin{questions}\item Montrer que $\sqrt2$ n’est pas un nombre entier.\item On suppose que $\sqrt2$ s’écrit avec un nombre fini de chiffres après la virgule, le dernier étant un chiffre $c$ non nul. Quel peut être le dernier chiffre de son carré ?\item Conclure.\end{questions}\label{fig-defi-4}\figureGeo{defi-racine}
\par\vspace{5mm}{\otLight\fontsize{8}{10}\selectfont\color{otMuted}Espace de recherche}\par\par\vspace{7mm}{\color{otRule!65}\hrule height .25pt}\par\vspace{7mm}{\color{otRule!65}\hrule height .25pt}\par\vspace{7mm}{\color{otRule!65}\hrule height .25pt}\par\vspace{7mm}{\color{otRule!65}\hrule height .25pt}\par\vspace{7mm}{\color{otRule!65}\hrule height .25pt}\par\vspace{7mm}{\color{otRule!65}\hrule height .25pt}
\end{blocEdition}
\editionPage{defi-ex-5}{D05 · Défi}{Comparer sans choisir de valeur}{Approfondissement · 30 à 60 minutes indicatives}
\begin{blocEdition}{defi-ex-5}
Pour $x>0$, un rectangle A mesure 3 cm sur $x+4$ cm. Un rectangle B mesure 2 cm sur $x+6$ cm. Quelle aire est la plus grande ? Est-il possible qu’elles soient égales pour une longueur $x$ strictement positive ?\label{fig-defi-5}\figureGeo{defi-litteral}
\par\vspace{5mm}{\otLight\fontsize{8}{10}\selectfont\color{otMuted}Espace de recherche}\par\par\vspace{7mm}{\color{otRule!65}\hrule height .25pt}\par\vspace{7mm}{\color{otRule!65}\hrule height .25pt}\par\vspace{7mm}{\color{otRule!65}\hrule height .25pt}\par\vspace{7mm}{\color{otRule!65}\hrule height .25pt}\par\vspace{7mm}{\color{otRule!65}\hrule height .25pt}\par\vspace{7mm}{\color{otRule!65}\hrule height .25pt}
\end{blocEdition}
\editionPage{defi-ex-6}{D06 · Défi}{Compter les pièces}{Approfondissement · 30 à 60 minutes indicatives}
\begin{blocEdition}{defi-ex-6}
Une boîte contient 42 pièces, uniquement de 1 euro et de 2 euros. Leur valeur totale est 68 euros. Trouver le nombre de pièces de chaque sorte et vérifier les deux contraintes.\par\vspace{5mm}{\otLight\fontsize{8}{10}\selectfont\color{otMuted}Espace de recherche}\par\par\vspace{7mm}{\color{otRule!65}\hrule height .25pt}\par\vspace{7mm}{\color{otRule!65}\hrule height .25pt}\par\vspace{7mm}{\color{otRule!65}\hrule height .25pt}\par\vspace{7mm}{\color{otRule!65}\hrule height .25pt}\par\vspace{7mm}{\color{otRule!65}\hrule height .25pt}\par\vspace{7mm}{\color{otRule!65}\hrule height .25pt}
\end{blocEdition}
\editionPage{defi-ex-7}{D07 · Défi}{Partager un bénéfice}{Approfondissement · 30 à 60 minutes indicatives}
\begin{blocEdition}{defi-ex-7}
Trois associés ont investi 12 000 €, 18 000 € et 30 000 €. Le bénéfice de 9 000 € est partagé proportionnellement aux sommes investies.\begin{questions}\item Écrire le ratio des investissements sous forme simplifiée et calculer les trois parts.\item L’année suivante, le troisième associé a retiré 10 000 €. Le bénéfice, de 10 000 €, est partagé de la même manière. Calculer les nouvelles parts.\item Quel associé voit sa part augmenter le plus ?\end{questions}\par\vspace{5mm}{\otLight\fontsize{8}{10}\selectfont\color{otMuted}Espace de recherche}\par\par\vspace{7mm}{\color{otRule!65}\hrule height .25pt}\par\vspace{7mm}{\color{otRule!65}\hrule height .25pt}\par\vspace{7mm}{\color{otRule!65}\hrule height .25pt}\par\vspace{7mm}{\color{otRule!65}\hrule height .25pt}\par\vspace{7mm}{\color{otRule!65}\hrule height .25pt}\par\vspace{7mm}{\color{otRule!65}\hrule height .25pt}
\end{blocEdition}
\editionPage{defi-ex-8}{D08 · Défi}{Annuler une remise}{Approfondissement · 30 à 60 minutes indicatives}
\begin{blocEdition}{defi-ex-8}
Un prix baisse de 20\,\%. De quel pourcentage faut-il augmenter le nouveau prix pour retrouver l’ancien ? Prouver que ce pourcentage ne dépend pas du prix initial.\par\vspace{5mm}{\otLight\fontsize{8}{10}\selectfont\color{otMuted}Espace de recherche}\par\par\vspace{7mm}{\color{otRule!65}\hrule height .25pt}\par\vspace{7mm}{\color{otRule!65}\hrule height .25pt}\par\vspace{7mm}{\color{otRule!65}\hrule height .25pt}\par\vspace{7mm}{\color{otRule!65}\hrule height .25pt}\par\vspace{7mm}{\color{otRule!65}\hrule height .25pt}\par\vspace{7mm}{\color{otRule!65}\hrule height .25pt}
\end{blocEdition}
\editionPage{defi-ex-9}{D09 · Défi}{Un découpage variable}{Approfondissement · 30 à 60 minutes indicatives}
\begin{blocEdition}{defi-ex-9}
Dans un rectangle de 18 cm sur 10 cm, on retire quatre carrés de côté $x$ cm aux coins, avec $0<x<5$. Donner la formule de l’aire restante $A(x)$. Parmi $x=2$, 3 et 4, lequel laisse 116 cm$^2$ ? Peut-on retrouver ce choix sans tester les trois valeurs ?\label{fig-defi-9}\figureGeo{coins-retires}
\par\vspace{5mm}{\otLight\fontsize{8}{10}\selectfont\color{otMuted}Espace de recherche}\par\par\vspace{7mm}{\color{otRule!65}\hrule height .25pt}\par\vspace{7mm}{\color{otRule!65}\hrule height .25pt}\par\vspace{7mm}{\color{otRule!65}\hrule height .25pt}\par\vspace{7mm}{\color{otRule!65}\hrule height .25pt}\par\vspace{7mm}{\color{otRule!65}\hrule height .25pt}\par\vspace{7mm}{\color{otRule!65}\hrule height .25pt}
\end{blocEdition}
\editionPage{defi-ex-10}{D10 · Défi}{Même moyenne, autre médiane}{Approfondissement · 30 à 60 minutes indicatives}
\begin{blocEdition}{defi-ex-10}
Comparer les séries A : 4, 4, 4, 16, 16, 16 et B : 0, 8, 8, 8, 8, 28. Calculer la moyenne, une médiane et l’étendue de chacune. Que peut-on conclure sur l’information fournie par une moyenne seule ?\par\vspace{5mm}{\otLight\fontsize{8}{10}\selectfont\color{otMuted}Espace de recherche}\par\par\vspace{7mm}{\color{otRule!65}\hrule height .25pt}\par\vspace{7mm}{\color{otRule!65}\hrule height .25pt}\par\vspace{7mm}{\color{otRule!65}\hrule height .25pt}\par\vspace{7mm}{\color{otRule!65}\hrule height .25pt}\par\vspace{7mm}{\color{otRule!65}\hrule height .25pt}\par\vspace{7mm}{\color{otRule!65}\hrule height .25pt}
\end{blocEdition}
\editionPage{defi-ex-11}{D11 · Défi}{Un jeu de dés équitable ?}{Approfondissement · 30 à 60 minutes indicatives}
\begin{blocEdition}{defi-ex-11}
On lance un dé rouge et un dé bleu, équilibrés. Léa gagne si le produit des deux nombres est pair ; sinon, c’est Hugo qui gagne.\begin{questions}\item Le jeu est-il équitable ?\item On change la règle : Léa gagne si la somme des deux nombres est paire. Le jeu devient-il équitable ?\end{questions}\label{fig-defi-11}\figureGeo{defi-des}
\par\vspace{5mm}{\otLight\fontsize{8}{10}\selectfont\color{otMuted}Espace de recherche}\par\par\vspace{7mm}{\color{otRule!65}\hrule height .25pt}\par\vspace{7mm}{\color{otRule!65}\hrule height .25pt}\par\vspace{7mm}{\color{otRule!65}\hrule height .25pt}\par\vspace{7mm}{\color{otRule!65}\hrule height .25pt}\par\vspace{7mm}{\color{otRule!65}\hrule height .25pt}\par\vspace{7mm}{\color{otRule!65}\hrule height .25pt}
\end{blocEdition}
\editionPage{defi-ex-12}{D12 · Défi}{Six triangles de même aire}{Approfondissement · 30 à 60 minutes indicatives}
\begin{blocEdition}{defi-ex-12}
Dans le triangle ABC, I, J et K sont les milieux de $[BC]$, $[CA]$ et $[AB]$. Les médianes $(AI)$, $(BJ)$ et $(CK)$ se coupent en G. Elles découpent le triangle en six petits triangles. Démontrer que ces six triangles ont la même aire.\label{fig-defi-12}\figureGeo{defi-medianes}
\par\vspace{5mm}{\otLight\fontsize{8}{10}\selectfont\color{otMuted}Espace de recherche}\par\par\vspace{7mm}{\color{otRule!65}\hrule height .25pt}\par\vspace{7mm}{\color{otRule!65}\hrule height .25pt}\par\vspace{7mm}{\color{otRule!65}\hrule height .25pt}\par\vspace{7mm}{\color{otRule!65}\hrule height .25pt}\par\vspace{7mm}{\color{otRule!65}\hrule height .25pt}\par\vspace{7mm}{\color{otRule!65}\hrule height .25pt}
\end{blocEdition}
\editionPage{defi-ex-13}{D13 · Défi}{Le pont le plus court}{Approfondissement · 30 à 60 minutes indicatives}
\begin{blocEdition}{defi-ex-13}
Une rivière a deux rives parallèles distantes de 1 km. Dans un repère d’unité 1 km, les rives sont les droites d’équations $y=2$ et $y=3$. Le village A est en $(0\,;0)$ et le village B en $(6\,;5)$. On construit un pont MN perpendiculaire aux rives, avec M sur la rive $y=2$. Où placer M pour que le trajet $AM+MN+NB$ soit le plus court ?\label{fig-defi-13}\figureGeo{defi-pont}
\par\vspace{5mm}{\otLight\fontsize{8}{10}\selectfont\color{otMuted}Espace de recherche}\par\par\vspace{7mm}{\color{otRule!65}\hrule height .25pt}\par\vspace{7mm}{\color{otRule!65}\hrule height .25pt}\par\vspace{7mm}{\color{otRule!65}\hrule height .25pt}\par\vspace{7mm}{\color{otRule!65}\hrule height .25pt}\par\vspace{7mm}{\color{otRule!65}\hrule height .25pt}\par\vspace{7mm}{\color{otRule!65}\hrule height .25pt}
\end{blocEdition}
\editionPage{defi-ex-14}{D14 · Défi}{Varignon et les aires}{Approfondissement · 30 à 60 minutes indicatives}
\begin{blocEdition}{defi-ex-14}
ABCD est un quadrilatère non croisé dont les diagonales se coupent à l’intérieur. I, J, K et L sont les milieux de $[AB]$, $[BC]$, $[CD]$ et $[DA]$. On sait que IJKL est un parallélogramme : c’est le théorème du mathématicien français Pierre Varignon. Démontrer que son aire est la moitié de celle de ABCD.\label{fig-defi-14}\figureGeo{defi-varignon}
\par\vspace{5mm}{\otLight\fontsize{8}{10}\selectfont\color{otMuted}Espace de recherche}\par\par\vspace{7mm}{\color{otRule!65}\hrule height .25pt}\par\vspace{7mm}{\color{otRule!65}\hrule height .25pt}\par\vspace{7mm}{\color{otRule!65}\hrule height .25pt}\par\vspace{7mm}{\color{otRule!65}\hrule height .25pt}\par\vspace{7mm}{\color{otRule!65}\hrule height .25pt}\par\vspace{7mm}{\color{otRule!65}\hrule height .25pt}
\end{blocEdition}
\editionPage{defi-ex-15}{D15 · Défi}{Le câble et ses attaches}{Approfondissement · 30 à 60 minutes indicatives}
\begin{blocEdition}{defi-ex-15}
Un câble doit suivre l’hypoténuse d’un triangle rectangle dont les côtés de l’angle droit mesurent 9 m et 12 m. Il faut en plus 0,75 m à chaque extrémité pour les attaches. Le câble coûte 4,80 euros le mètre. Calculer la longueur et le coût exacts.\label{fig-defi-15}\figureGeo{defi-pyth}
\par\vspace{5mm}{\otLight\fontsize{8}{10}\selectfont\color{otMuted}Espace de recherche}\par\par\vspace{7mm}{\color{otRule!65}\hrule height .25pt}\par\vspace{7mm}{\color{otRule!65}\hrule height .25pt}\par\vspace{7mm}{\color{otRule!65}\hrule height .25pt}\par\vspace{7mm}{\color{otRule!65}\hrule height .25pt}\par\vspace{7mm}{\color{otRule!65}\hrule height .25pt}\par\vspace{7mm}{\color{otRule!65}\hrule height .25pt}
\end{blocEdition}
\editionPage{defi-ex-16}{D16 · Défi}{Le plus grand triangle du demi-cercle}{Approfondissement · 30 à 60 minutes indicatives}
\begin{blocEdition}{defi-ex-16}
$[AB]$ est un diamètre de 10 cm d’un cercle de centre O, et M un point du cercle distinct de A et B.\begin{questions}\item Démontrer que $AM^2+BM^2=100$, quelle que soit la position de M.\item Démontrer que l’aire du triangle AMB est au plus égale à 25 cm$^2$. Pour quelle position de M est-elle égale à 25 cm$^2$ ?\end{questions}\label{fig-defi-16}\figureGeo{defi-cercle}
\par\vspace{5mm}{\otLight\fontsize{8}{10}\selectfont\color{otMuted}Espace de recherche}\par\par\vspace{7mm}{\color{otRule!65}\hrule height .25pt}\par\vspace{7mm}{\color{otRule!65}\hrule height .25pt}\par\vspace{7mm}{\color{otRule!65}\hrule height .25pt}\par\vspace{7mm}{\color{otRule!65}\hrule height .25pt}\par\vspace{7mm}{\color{otRule!65}\hrule height .25pt}\par\vspace{7mm}{\color{otRule!65}\hrule height .25pt}
\end{blocEdition}
\editionPage{defi-ex-17}{D17 · Défi}{Une pyramide de verre}{Approfondissement · 30 à 60 minutes indicatives}
\begin{blocEdition}{defi-ex-17}
On modélise la pyramide du Louvre par une pyramide régulière à base carrée de 35 m de côté et de 21 m de hauteur, valeurs arrondies. O est le centre de la base et M le milieu d’un côté de la base.\begin{questions}\item Calculer le volume de cette pyramide.\item Calculer SM, la hauteur d’une face triangulaire, au centième de mètre.\item En déduire l’aire des quatre faces vitrées, au mètre carré près.\end{questions}\label{fig-defi-17}\figureGeo{defi-pyramide}
\par\vspace{5mm}{\otLight\fontsize{8}{10}\selectfont\color{otMuted}Espace de recherche}\par\par\vspace{7mm}{\color{otRule!65}\hrule height .25pt}\par\vspace{7mm}{\color{otRule!65}\hrule height .25pt}\par\vspace{7mm}{\color{otRule!65}\hrule height .25pt}\par\vspace{7mm}{\color{otRule!65}\hrule height .25pt}\par\vspace{7mm}{\color{otRule!65}\hrule height .25pt}\par\vspace{7mm}{\color{otRule!65}\hrule height .25pt}
\end{blocEdition}
\editionPage{defi-ex-18}{D18 · Défi}{La suite de Syracuse}{Approfondissement · 30 à 60 minutes indicatives}
\begin{blocEdition}{defi-ex-18}
Le script demande un entier n puis répète 16 fois : si le reste de n divisé par 2 vaut 0, mettre n à n divisé par 2 ; sinon, mettre n à $3\times n+1$ ; dire n.\begin{questions}\item Donner les seize valeurs dites pour l’entrée 7.\item À quel tour obtient-on 1 pour la première fois ? Que se passerait-il si l’on continuait ?\item Le script n’utilise qu’une variable. Expliquer pourquoi c’est suffisant.\end{questions}\label{fig-defi-18}\figureGeo{defi-script}
\par\vspace{5mm}{\otLight\fontsize{8}{10}\selectfont\color{otMuted}Espace de recherche}\par\par\vspace{7mm}{\color{otRule!65}\hrule height .25pt}\par\vspace{7mm}{\color{otRule!65}\hrule height .25pt}\par\vspace{7mm}{\color{otRule!65}\hrule height .25pt}\par\vspace{7mm}{\color{otRule!65}\hrule height .25pt}\par\vspace{7mm}{\color{otRule!65}\hrule height .25pt}\par\vspace{7mm}{\color{otRule!65}\hrule height .25pt}
\end{blocEdition}
